% chapters/02_grundlagen.tex
% Zweites Kapitel: Theoretische Grundlagen
% Dieses Kapitel zeigt, wie Zitate, Tabellen und Labels funktionieren.

\chapter{Theoretische Grundlagen}
\label{chap:grundlagen}

Hier steht der einleitende Text des Kapitels \autocite{beispiel_buch_2020}.
Fuer mehrere Quellen gleichzeitig: \autocite{beispiel_artikel_2019, beispiel_online_2024}.

\section{Erster Abschnitt}
\label{sec:abschnitt1}

Text des ersten Abschnitts. Ein Zitat im Fliestext: Laut
\textcite{beispiel_buch_2020} gilt \ldots

Mit Seitenangabe: \autocite[S.~42]{beispiel_buch_2020}.

\section{Tabellen}
\label{sec:tabellen}

Tabelle~\ref{tab:beispiel} zeigt ein Beispiel fuer eine professionell
formatierte Tabelle mit dem \texttt{booktabs}-Paket.

\begin{table}[h]
\centering
\caption{Beispieltabelle}
\label{tab:beispiel}
\begin{tabular}{lcc}
\toprule
\textbf{Kriterium} & \textbf{Option A} & \textbf{Option B} \\
\midrule
Merkmal 1 & Gut    & Mittel \\
Merkmal 2 & Mittel & Gut    \\
Merkmal 3 & Gut    & Gut    \\
\bottomrule
\end{tabular}
\end{table}

\section{Abbildungen}
\label{sec:abbildungen}

% Abbildung einbinden (Datei im figures/-Ordner ablegen):
% \begin{figure}[h]
%   \centering
%   \includegraphics[width=0.8\textwidth]{figures/mein-bild.png}
%   \caption{Beschreibung der Abbildung}
%   \label{fig:mein-bild}
% \end{figure}
